\section{Benchmarks and Experimental Settings}
\label{sec:xp_settings}

We evaluate contextual privacy using two settings. The \textbf{probing} setting uses targeted, single-turn queries to efficiently test a model's \textit{explicit} privacy understanding. The \textbf{agentic} setting simulates multi-turn interactions in real web environments to assess \textit{implicit} privacy understanding, with greater complexity and cost. As recommended by \citet{privacylens}, we use both settings to ensure a comprehensive assessment of utility--privacy trade-offs. 

\paragraph{Probing setting.}

Our probing evaluation uses \textbf{AirGapAgent-R}, a reconstruction of the unavailable AirGapAgent benchmark~\citep{bagdasarian2024airgapagent}, based on the authors' public methodology (procedure in Appendix~\ref{app:airgapagent}). The dataset includes 20 synthetic user profiles, each with 26 data fields (e.g., name, age, health conditions), evaluated across 8 scenarios (e.g., restaurant or medical bookings), totaling 4{,}160 datapoints.
Each prompt presents the model with a user profile, a scenario, and a question about whether a specific data field should be shared. Ground-truth labels for each prompt indicate when sharing is appropriate (e.g., age for a doctor’s appointment) or not.
We report a utility score as the percentage of appropriate fields correctly shared, and a privacy score as the percentage of examples where no inappropriate fields are leaked (higher is better for both). Sensitive data is detected using a \texttt{gpt-4o-mini}-based extractor applied to both the final answer and the reasoning trace. We will release AirGapAgent-R to support future research on contextual privacy.

\paragraph{Agentic setting.}
We use the AgentDAM benchmark \cite{agentdam} to evaluate contextual privacy in simulated web environments, split across three domains: shopping, Reddit, and GitLab. Models interact with websites via a textual accessibility tree, contextual input (e.g., user chat), and a set of predefined actions. Agents act step-by-step until issuing the \texttt{stop} action or reaching 10 actions. We report a utility score (task success rate) and a privacy score (proportion of tasks without leakage in any interaction). Privacy is assessed for both answer and reasoning using \texttt{gpt-4o-mini} with a four-shot prompt, following the original setup (prompts in Appendix~\ref{app:prompts}).

\paragraph{Models evaluated and prompting techniques.} 

We evaluate 13 models ranging from 8B to over 600B parameters, grouped by family to reflect shared lineage through distillation. We compare vanilla LLMs, CoT-prompted vanilla models, and Large Reasoning Models. Distilled models (e.g., DeepSeek’s \texttt{R1-} variants of \texttt{Llama} and \texttt{Qwen}) are included, alongside others such as \texttt{QwQ}, \texttt{s1}, and \texttt{s1.1}. In probing, we ask the model to maintain thinking within \texttt{<think>} and \texttt{</think>} and to anonymize sensitive data in the reasoning using placeholders (e.g., \texttt{<name>}); in the agentic setup, we apply the CoT mitigation from AgentDAM. Model specifications and configuration details, along with complete prompt templates (including both system and evaluator prompts), are provided in Appendix~\ref{app:hparams} and \ref{app:prompts}.
Results are averaged over seeds (probing) or splits (agentic), with metric variation reported in percentage points (\%p.).


